\chapter{Invariant Laws, Four-Vectors, and the Field Tensor}

We shall begin with a law that nature could have obeyed but does not. Its
failure will tell us how a physical law ought to be written. The same test
first leads from coordinate components to ordinary vectors, then from
ordinary vectors to four-vectors, and finally from four-vectors to tensors.
By the end, the electric and magnetic fields will no longer be separate
objects: together they will occupy the six independent slots of one
antisymmetric tensor.

Unless stated otherwise, we use units in which \(c=1\), order spacetime
components as \((1,2,3,0)=(x,y,z,t)\), and use the metric signature
\((-,-,-,+)\).

\section{A possible law with the wrong symmetry}

Fix an \(x\)- and a \(y\)-axis permanently in space. Now decree that two
particles interact whenever their \(x\)-coordinates agree. The interaction
might reverse their motion or send them away in new directions; its detailed
dynamics does not matter. The proposed trigger is simply
\begin{equation}
  x_1-x_2=0,
\end{equation}
where the subscripts label particles rather than coordinate components.
There is no mathematical contradiction in this rule. The question is what
makes it physically implausible.

\begin{classroomqa}
\audiencequestion{Does the rule violate the homogeneity of space?}
\lecturerresponse{No. Translating the origin changes both coordinates by the
same amount and leaves \(x_1-x_2\) unchanged. The rule specifies no preferred
place. Its defect is isotropy: it singles out one preferred direction.}
\end{classroomqa}

Rotate the axes through \(90^\circ\). The old \(x\)-direction is now called
\(y\), so the same physical events would be described by equality of the
\(y\)-coordinates. At a general angle the condition would involve a linear
combination of \(x\) and \(y\). The form of the law therefore depends on how
the blackboard was oriented. It is translation invariant but not rotation
invariant.

We usually give physical laws their simplest form in orthogonal coordinates.
One may use more general coordinates, but then the metric and the formulas
become less simple. Within the class of orthogonal Cartesian frames, no
orientation should be privileged.

The repair is immediate. In two dimensions require both coordinates to
agree; in three dimensions require all three:
\begin{equation}
  x_1=x_2,
  \qquad y_1=y_2,
  \qquad z_1=z_2.
\end{equation}
The particles can then interact only when they reach the same point, even if
they approach it along different lines.

\begin{classroomqa}
\audiencequestion{Must the \(z\)-coordinates and the times agree as well?}
\lecturerresponse{In three-dimensional space the \(z\)-coordinates must also
agree. Time is being suppressed for the moment; it will be restored as soon
as the spatial argument is complete.}
\end{classroomqa}

Let \(\boldsymbol r_a\) be the position of particle \(a\). The three component
conditions become one vector equation,
\begin{equation}
  \boldsymbol r_1-\boldsymbol r_2=\boldsymbol 0.
\end{equation}
A component can vanish in one rotated frame and fail to vanish in another.
The whole vector is different: if every component is zero in one frame, every
component is zero in every rotated frame. Likewise, if two vectors are equal
in one frame, they remain equal after both are transformed.

\begin{figure}[H]
\centering
\includegraphics[width=0.78\textwidth]{lecture_06_figure_01.png}
\caption{The three coordinate equalities that make collision mean spatial
coincidence rather than coincidence along one preferred axis.}
\lectureframe{00:09:00}
\end{figure}

This is why vector equations are useful. Their components change, but the
object transforms coherently. A law written for the complete object can
retain its form when the coordinates are rotated.

\section{Coincidence in spacetime}

Now replace the vertical spatial axis by time. A Newtonian rule might say
that two changes of motion occur only at equal times. In Newtonian mechanics
simultaneity is shared by all inertial observers, so such a rule is possible.
In special relativity it selects a preferred frame. Events simultaneous in
one frame need not be simultaneous in another.

Demanding equality only of the spatial coordinates is no better. A
spacetime displacement can have vanishing spatial components and a nonzero
time component,
\begin{equation}
  \Delta X^\mu=(0,0,0,\Delta t),
\end{equation}
and a boost mixes \(\Delta t\) with a spatial component. The invariant local
statement is that the two changes of motion occur at the same event. Thus a
simple collision requires
\begin{equation}
  X_1^\mu=X_2^\mu,
  \qquad\text{or}\qquad
  X_1^\mu-X_2^\mu=0.
\end{equation}
All four coordinates agree: the particles occupy the same place at the same
time.

The collision example gives a general rule for writing physics. We may put
everything on one side and set a complete four-vector equal to zero, or set
one complete four-vector equal to another. Equating disconnected pieces of
a four-vector does not, by itself, give a Lorentz-invariant statement.

\section{Boosts, rotations, and matrix notation}

A four-vector is an object that transforms as a spacetime displacement does.
To make that definition useful, first recall the transformations of the
coordinates themselves. For a boost of speed \(v\) in the \(x\)-direction,
\begin{align}
  x'&=\gamma(x-vt), & y'&=y, & z'&=z,
  & t'&=\gamma(t-vx),                                      \label{eq:l6-boost}
\end{align}
where
\begin{equation}
  \gamma(v)=\frac{1}{\sqrt{1-v^2}}.
\end{equation}

\begin{classroomqa}
\audiencequestion{Has the speed of light been set equal to one in these
formulas?}
\lecturerresponse{Yes. With units restored, take \(X^0=ct\), replace
\(1-v^2\) by \(1-v^2/c^2\), and restore the corresponding factors of \(c\)
in the mixed space--time terms.}
\end{classroomqa}

A rotation about the \(z\)-axis mixes \(x\) and \(y\), leaving \(z\) and
\(t\) alone:
\begin{equation}
\begin{pmatrix}
  x'\\ y'\\ z'\\ t'
\end{pmatrix}
=
\begin{pmatrix}
  \cos\theta & \sin\theta & 0 & 0\\
  -\sin\theta & \cos\theta & 0 & 0\\
  0 & 0 & 1 & 0\\
  0 & 0 & 0 & 1
\end{pmatrix}
\begin{pmatrix}
  x\\ y\\ z\\ t
\end{pmatrix}.
\end{equation}
The zeros are worth displaying: they show explicitly which components do not
enter a given transformed coordinate. In the same ordering, the boost in
Eq.~\eqref{eq:l6-boost} is
\begin{equation}
\begin{pmatrix}
  x'\\ y'\\ z'\\ t'
\end{pmatrix}
=
\begin{pmatrix}
  \gamma & 0 & 0 & -\gamma v\\
  0 & 1 & 0 & 0\\
  0 & 0 & 1 & 0\\
  -\gamma v & 0 & 0 & \gamma
\end{pmatrix}
\begin{pmatrix}
  x\\ y\\ z\\ t
\end{pmatrix}.
\end{equation}

Rotations and boosts can be composed by matrix multiplication. To boost in a
direction not aligned with the original \(x\)-axis, one may rotate that
direction into the \(x\)-axis, apply the boost, and rotate as needed. Products
of these elementary matrices generate more general Lorentz transformations.

\begin{classroomqa}
\audiencequestion{Does it matter whether one rotates first and then boosts,
or boosts first and then rotates?}
\lecturerresponse{Yes. After the first operation the axes on which the second
one acts have changed. The two matrix products need not agree, so the order of
operations matters.}
\end{classroomqa}

All these linear transformations can be written at once as
\begin{equation}
  X'^{\mu}=\Lambda^{\mu}{}_{\nu}X^{\nu}.                  \label{eq:l6-vector-transform}
\end{equation}
The repeated index \(\nu\) is summed over \(1,2,3,0\). This is simply the
usual row-by-column sum in matrix multiplication, expressed in index
notation.

\begin{figure}[H]
\centering
\includegraphics[width=0.78\textwidth]{lecture_06_figure_02.png}
\caption{The indexed Lorentz transformation above the beginning of its
column-vector form.}
\lectureframe{00:31:55}
\end{figure}

Equation~\eqref{eq:l6-vector-transform} is also the definition of a general
contravariant four-vector. If \(A^\mu\) is such a vector, then
\begin{equation}
  A'^{\mu}=\Lambda^{\mu}{}_{\nu}A^{\nu}.
\end{equation}
Its four components need not represent position. What makes the object a
four-vector is its transformation law.

\section{Scalars, the metric, and invariant products}

Before using more four-vectors, it is helpful to step back to the simpler
notion of a scalar. Under spatial rotations, the temperature at a specified
point is a scalar, and so is the distance between two specified points. Each
is one number, independent of the orientation of the axes. One could even
make the deliberately silly numerical claim that a distance measured in
centimetres equals the room temperature measured in degrees Fahrenheit. The
claim has no physical depth, but because both sides are rotational scalars,
its truth or falsity does not depend on how the axes are turned.

The components of a vector are not scalars. We therefore need a systematic
way to combine them into scalars. Write a general contravariant four-vector as
\begin{equation}
  A^\mu=(A^1,A^2,A^3,A^0).
\end{equation}
The covariant components describe the same geometric object with a lower
index. With our signature they are obtained using the Minkowski metric,
\begin{equation}
  A_\mu=\eta_{\mu\nu}A^\nu,
  \qquad
  \eta_{\mu\nu}=
  \begin{pmatrix}
    -1&0&0&0\\
    0&-1&0&0\\
    0&0&-1&0\\
    0&0&0&1
  \end{pmatrix},                                          \label{eq:l6-metric}
\end{equation}
so that
\begin{equation}
  A_\mu=(-A^1,-A^2,-A^3,A^0).
\end{equation}

\begin{classroomqa}
\audiencequestion{Are two indices and signs on the board inconsistent?}
\lecturerresponse{Yes. The spatial component in the displayed column needs
the corresponding minus sign, and the components on the right of the
index-lowering equation carry upper indices before multiplication by
\(\eta_{\mu\nu}\). The time label is \(0\), not \(4\).}
\end{classroomqa}

\begin{figure}[H]
\centering
\includegraphics[width=0.8\textwidth]{lecture_06_figure_03.png}
\caption{Contravariant and covariant components, the metric matrix, and the
beginning of the invariant contraction.}
\lectureframe{00:45:00}
\end{figure}

Index lowering is not decoration. It packages the signs required for the
Lorentz-invariant norm:
\begin{align}
  A_\mu A^\mu
  &=-(A^1)^2-(A^2)^2-(A^3)^2+(A^0)^2,                    \label{eq:l6-norm}\\
  X_\mu X^\mu
  &=t^2-x^2-y^2-z^2.
\end{align}
The second expression is the invariant spacetime interval from the origin.
Because \(A^\mu\) transforms in the same way as \(X^\mu\), its contraction
has the same invariance.

The notation also teaches a useful bookkeeping rule. A repeated summation
index appears once upstairs and once downstairs. For example,
\begin{equation}
  A_\mu A^\mu
  =\eta_{\mu\nu}A^\mu A^\nu.
\end{equation}
The names of summed, or dummy, indices may be changed without changing the
quantity.

\begin{classroomqa}
\audiencequestion{Should one of these vectors be regarded as a row and the
other as a column, and does their order matter?}
\lecturerresponse{For the present calculation, regard the indexed quantities
as components and let the metric perform the contraction; there is no need to
turn the discussion into row-versus-column operator algebra. These are
ordinary numbers, so their order does not matter. Noncommuting operators
would require more care, but that is not the subject here.}
\end{classroomqa}

The same construction gives an inner product between two different
four-vectors,
\begin{equation}
  A\mathbin{\cdot}B=A_\mu B^\mu=A^\mu B_\mu.
\end{equation}
Its invariance follows without writing an explicit Lorentz matrix. Since
\(A+B\) and \(A-B\) are four-vectors, both of their squared norms are
invariant:
\begin{align}
  (A+B)^2&=A^2+B^2+2A\mathbin{\cdot}B,\\
  (A-B)^2&=A^2+B^2-2A\mathbin{\cdot}B.
\end{align}
Subtracting gives
\begin{equation}
  (A+B)^2-(A-B)^2=4A\mathbin{\cdot}B,                    \label{eq:l6-polarization}
\end{equation}
so \(A\mathbin{\cdot}B\) is invariant as well. The factor of four follows
directly from the expansion.\editorialnote{The board argument supplies this
polarization identity; writing out the factor makes the spoken cancellation
unambiguous.}

\begin{figure}[H]
\centering
\includegraphics[width=0.8\textwidth]{lecture_06_figure_04.png}
\caption{The two invariant squares used to isolate the inner product of
different four-vectors.}
\lectureframe{00:52:20}
\end{figure}

Invariant quantities are natural ingredients of physical laws: unlike an
isolated component, they carry the same value from frame to frame.

\section{Proper velocity and four-momentum}

Consider two neighboring events on a particle's worldline. Their separation
is the infinitesimal four-vector
\begin{equation}
  dX^\mu=(dx,dy,dz,dt).
\end{equation}
For a timelike worldline, proper time is defined by
\begin{equation}
  d\tau^2=dX_\mu dX^\mu
  =dt^2-dx^2-dy^2-dz^2.                                  \label{eq:l6-proper-time}
\end{equation}
The proper four-velocity is
\begin{equation}
  U^\mu=\frac{dX^\mu}{d\tau}.                            \label{eq:l6-four-velocity}
\end{equation}
It measures the rate of change of spacetime position with respect to the
particle's own proper time, not with respect to the coordinate time of a
chosen frame.

Although \(U^\mu\) has four components, only three are independent. Dividing
Eq.~\eqref{eq:l6-proper-time} by \(d\tau^2\) gives
\begin{equation}
  U_\mu U^\mu=1.
\end{equation}
Thus the four-velocity is a unit timelike vector. The future-directed choice
fixes the sign of its time component.

\begin{classroomqa}
\audiencequestion{At speeds much smaller than \(c\), does four-velocity reduce
to ordinary velocity?}
\lecturerresponse{Its three spatial components approach ordinary velocity,
while its time component approaches one in \(c=1\) units. The relation follows
by applying the chain rule and computing \(dt/d\tau\).}
\end{classroomqa}

For each spatial component,
\begin{equation}
  U^i=\frac{dx^i}{d\tau}
     =\frac{dx^i}{dt}\frac{dt}{d\tau}
     =v^i\frac{dt}{d\tau}.
\end{equation}
Divide Eq.~\eqref{eq:l6-proper-time} by \(dt^2\):
\begin{equation}
  \left(\frac{d\tau}{dt}\right)^2
  =1-v_x^2-v_y^2-v_z^2
  =1-\boldsymbol v^2.
\end{equation}
Therefore
\begin{equation}
  \frac{dt}{d\tau}=\frac{1}{\sqrt{1-\boldsymbol v^2}}
  =\gamma(\lvert\boldsymbol v\rvert),
\end{equation}
and
\begin{equation}
  U^\mu=(\gamma v_x,\gamma v_y,\gamma v_z,\gamma).
                                                                  \label{eq:l6-u-components}
\end{equation}

\begin{figure}[H]
\centering
\includegraphics[width=0.8\textwidth]{lecture_06_figure_05.png}
\caption{The proper-time calculation that produces \(dt/d\tau=\gamma\) and
relates proper velocity to ordinary velocity.}
\lectureframe{01:02:00}
\end{figure}

\begin{classroomqa}
\audiencequestion{Does \(dx^2/dt^2\) mean a second derivative, and should
squaring a vector produce cross terms?}
\lecturerresponse{Here it means \((dx/dt)^2\), the square of a first
derivative. Moreover, \(\boldsymbol v^2\) means the Euclidean dot product
\(v_x^2+v_y^2+v_z^2\), not the square of the formal sum
\(v_x+v_y+v_z\); there are no cross terms.}
\end{classroomqa}

\begin{classroomqa}
\audiencequestion{Is the \(v\) in \(\gamma(v)\) the relative speed of two
coordinate frames or the speed of the particle?}
\lecturerresponse{Here it is the particle's speed in the chosen frame. The
function \(\gamma(v)=1/\sqrt{1-v^2}\) can be evaluated for any speed: in a
boost it is evaluated at the relative frame speed; in four-velocity it is
evaluated at the particle speed.}
\end{classroomqa}

For \(v\ll1\), \(\gamma\simeq1\), so the spatial components of \(U^\mu\)
approach \(\boldsymbol v\) and \(U^0\) approaches one. As the ordinary speed
approaches the speed of light, \(\gamma\) and all nonzero components of the
four-velocity grow without bound.

\begin{classroomqa}
\audiencequestion{Should not \(U^0\) have the units of a velocity rather than
being close to the number one?}
\lecturerresponse{The units disappeared when \(c\) was set to one. With
\(X^0=ct\), one has
\(U^\mu=(\gamma\boldsymbol v,\gamma c)\), so the time component approaches
\(c\) and \(U_\mu U^\mu=c^2\).}
\end{classroomqa}

For a particle of rest mass \(m\), define four-momentum by
\begin{equation}
  P^\mu=mU^\mu.
\end{equation}
In natural units,
\begin{equation}
  P^\mu=(\gamma m\boldsymbol v,\gamma m).
\end{equation}
The first three components approach the Newtonian momentum \(m\boldsymbol v\)
at low speed. Restoring units gives
\begin{equation}
  P^\mu=(\gamma m\boldsymbol v,E/c),
  \qquad E=\gamma mc^2,
\end{equation}
and hence
\begin{equation}
  E=mc^2+\frac{1}{2}m\boldsymbol v^2+O(mv^4/c^2).
\end{equation}
The time component is energy in \(c=1\) units, or energy divided by \(c\) when
all four components are assigned momentum units.

\begin{classroomqa}
\audiencequestion{Is ordinary \(\boldsymbol v\), with a fourth entry appended,
itself a four-vector?}
\lecturerresponse{No. Ordinary three-velocity does not transform as a
four-vector. The factor \(\gamma\) and the proper-time definition are
essential; the four-vector is \(U^\mu\), not \(\boldsymbol v\).}
\end{classroomqa}

\section{Why electric and magnetic fields must mix}

We now take what was called a short vacation from relativity. Start with the
pre-Einstein equation of motion and ask how it must be reorganized before it
can be valid in every inertial frame. Ordinary acceleration is the time
derivative of ordinary velocity. Its four-dimensional analogue is naturally
\begin{equation}
  A^\mu=\frac{dU^\mu}{d\tau}.
\end{equation}
At low speed its spatial components approach ordinary acceleration. The full
relativistic equation will be postponed; first inspect the old force law.

For a charge \(q\) in electric and magnetic fields,
\begin{equation}
  \boldsymbol F=m\boldsymbol a
  =q\left(\boldsymbol E+\boldsymbol v\times\boldsymbol B\right).
                                                                  \label{eq:l6-lorentz-force}
\end{equation}
The electric term is independent of the particle velocity, whereas the
magnetic term is not.

\begin{figure}[H]
\centering
\includegraphics[width=0.8\textwidth]{lecture_06_figure_06.png}
\caption{The pre-relativistic form of the Lorentz force law that motivates
combining \(\boldsymbol E\) and \(\boldsymbol B\).}
\lectureframe{01:18:20}
\end{figure}

Suppose an observer sees no electric field but does see a magnetic field and
a moving charge. The charge feels a magnetic force. In the instantaneous rest
frame of the charge, however, \(\boldsymbol v=0\), so the magnetic term in
Eq.~\eqref{eq:l6-lorentz-force} vanishes. The physical worldline has not
become inertial merely because the coordinates changed; the force in the new
frame must therefore include an electric contribution. More carefully, a
nonzero four-acceleration cannot be transformed to zero by changing inertial
frames, even though the three-acceleration and its components do transform.
\editorialnote{This four-acceleration statement is the invariant form of the
classroom argument and avoids treating ordinary three-acceleration itself as
frame invariant.}

Electric and magnetic fields must consequently transform into one another.
Together they have six components, so they are neither one scalar nor one
four-vector. We need a new transforming object.

\section{Tensors: one transformation per index}

A scalar may be called a tensor of rank zero; it has no indices. A four-vector
is a tensor of rank one; it has one index with four possible values. The
simplest rank-two example is a product
\begin{equation}
  A^\mu B^\nu.
\end{equation}
Each index has four values, so the object has \(4^2=16\) components, which may
be displayed as a \(4\times4\) array. A product of three vectors,
\begin{equation}
  A^\mu B^\nu C^\sigma,
\end{equation}
has \(4^3=64\) components and is a rank-three tensor. Higher ranks continue
the pattern. The name is less important than the rule that defines the
object.

If
\begin{equation}
  A'^{\mu}=\Lambda^\mu{}_{\nu}A^\nu,
  \qquad
  B'^{\sigma}=\Lambda^\sigma{}_{\tau}B^\tau,
\end{equation}
then their product transforms as
\begin{equation}
  A'^{\mu}B'^{\sigma}
  =\Lambda^\mu{}_{\nu}\Lambda^\sigma{}_{\tau}
   A^\nu B^\tau.
\end{equation}
This gives the defining law for a contravariant rank-two tensor:
\begin{equation}
  T'^{\mu\sigma}
  =\Lambda^\mu{}_{\nu}\Lambda^\sigma{}_{\tau}T^{\nu\tau}.
                                                                  \label{eq:l6-tensor-transform}
\end{equation}
There is one Lorentz matrix for each index. The repeated \(\nu\) and \(\tau\)
are summed; the open \(\mu\) and \(\sigma\) label the component being
calculated in the primed frame.

\begin{figure}[H]
\centering
\includegraphics[width=0.82\textwidth]{lecture_06_figure_07.png}
\caption{The transformation of two vectors and the resulting two-matrix
transformation law for a rank-two tensor.}
\lectureframe{01:27:15}
\end{figure}

A tensor need not be a single product of vectors. Both
\(A^\mu B^\nu\) and \(A^\nu B^\mu\) obey the tensor transformation law, as do
their sum and any sum such as
\begin{equation}
  A^\mu B^\nu+C^\mu D^\nu.
\end{equation}
The transformation law, not a particular factorization, is the definition.
A third-rank tensor acquires a third Lorentz matrix, and so on.

\section{Symmetric and antisymmetric tensors}

Two classes of rank-two tensors will be especially useful. A symmetric tensor
satisfies
\begin{equation}
  T^{\mu\nu}=T^{\nu\mu},
\end{equation}
so entries reflected across the diagonal agree. Its four diagonal entries
and six entries on one side of the diagonal determine the rest, giving
\begin{equation}
  4+6=10
\end{equation}
independent components.

An antisymmetric tensor satisfies
\begin{equation}
  T^{\mu\nu}=-T^{\nu\mu}.                                \label{eq:l6-antisymmetry}
\end{equation}
Putting \(\mu=\nu\) in Eq.~\eqref{eq:l6-antisymmetry} gives
\(T^{\mu\mu}=-T^{\mu\mu}\), so every diagonal entry vanishes. The six
entries in one triangular half determine the other six by a sign change.
There are therefore exactly six independent components.

\begin{figure}[H]
\centering
\includegraphics[width=0.82\textwidth]{lecture_06_figure_08.png}
\caption{A completed antisymmetric \(4\times4\) array: zeros on the diagonal
and opposite signs across it.}
\lectureframe{01:34:25}
\end{figure}

That count is the clue. The three electric components and three magnetic
components can fit into one antisymmetric rank-two tensor.

\section{The electromagnetic field tensor}

Let the spacetime labels run across rows and columns in the order
\((1,2,3,0)=(x,y,z,t)\). The arrangement written on the board is
\begin{equation}
  F_{\mu\nu}=
  \begin{pmatrix}
    0    & -B_3 & B_2  & E_1\\
    B_3  & 0    & -B_1 & E_2\\
    -B_2 & B_1  & 0    & E_3\\
    -E_1 & -E_2 & -E_3 & 0
  \end{pmatrix}.                                         \label{eq:l6-field-tensor}
\end{equation}
The purely spatial slots contain the magnetic field; the slots carrying one
spatial and one time index contain the electric field. Antisymmetry supplies
the lower triangular entries once the upper triangular entries are known.

\begin{classroomqa}
\audiencequestion{Are the magnetic entries all negative, and do the labels
correspond to \(x,y,z\)?}
\lecturerresponse{The signs in the upper spatial triangle alternate as shown:
\(-B_3,+B_2,-B_1\). The labels \(1,2,3\) correspond to \(x,y,z\), while
\(0\) labels time. Reflection across the diagonal reverses every sign.}
\end{classroomqa}

\begin{figure}[H]
\centering
\includegraphics[width=0.82\textwidth]{lecture_06_figure_09.png}
\caption{The completed electromagnetic field tensor in the lecture's
spatial-first component ordering.}
\lectureframe{01:41:05}
\end{figure}

The notation is conventionally \(F_{\mu\nu}\), rather than a generic
\(T_{\mu\nu}\). The lecturer conjectures that the letter might refer to
Faraday, while explicitly saying that he does not know the history of the
notation. What matters here is its antisymmetry and its six field components,
not that conjecture about the letter.

\begin{classroomqa}
\audiencequestion{Does each matrix entry represent the strength of a field in
the indicated coordinate plane?}
\lecturerresponse{It is a component of the electromagnetic field tensor.
The entries involving time are electric components, and the spatial--spatial
entries encode magnetic components.}
\end{classroomqa}

Fence off the row and column carrying the time index. The remaining
\(3\times3\) antisymmetric block is magnetic. Its independent entries may be
remembered cyclically:
\begin{equation}
  F_{12}=-B_3,
  \qquad F_{23}=-B_1,
  \qquad F_{31}=-B_2.                                   \label{eq:l6-magnetic-cycle}
\end{equation}
Starting from the first relation, advance every label around
\(1\to2\to3\to1\). Equivalently,
\begin{equation}
  F_{ij}=-\epsilon_{ijk}B_k,
\end{equation}
with \(\epsilon_{123}=+1\).\editorialnote{The Levi--Civita formula is the
compact standard form of the cyclic rule developed explicitly on the board.}

\begin{figure}[H]
\centering
\includegraphics[width=0.82\textwidth]{lecture_06_figure_10.png}
\caption{The cyclic spatial components that encode the magnetic field in the
antisymmetric tensor.}
\lectureframe{01:46:20}
\end{figure}

The electric identifications are more direct:
\begin{equation}
  E_1=F_{10},
  \qquad E_2=F_{20},
  \qquad E_3=F_{30}.
\end{equation}

The magnetic field consequently has two useful identities. Within
three-dimensional space, a three-component axial vector can be paired with a
\(3\times3\) antisymmetric tensor. In spacetime, however, it is only the
spatial part of the six-component tensor \(F_{\mu\nu}\). This pairing is
special to three dimensions: an antisymmetric \(4\times4\) tensor has six
independent components and cannot be identified with a four-vector.
\editorialnote{Under proper rotations the magnetic field behaves as a vector;
under reflections it is more precisely an axial vector, or pseudovector.}

The purpose of this construction is not component counting for its own sake.
Because \(F_{\mu\nu}\) transforms as a tensor, a Lorentz boost automatically
mixes its electric and magnetic entries. In the next lecture it can be
combined with four-velocity and four-acceleration to replace
Eq.~\eqref{eq:l6-lorentz-force} by a covariant equation of motion.

\section{Questions about Einstein and Minkowski}

The final minutes turn from the formalism to how it arose historically. The
answers are offered as recollections rather than a settled history, and that
qualification is part of the exchange.

\begin{classroomqa}
\audiencequestion{Was Einstein seeking a Lorentz-invariant formulation on
conceptual grounds, or was he driven by particular observations?}
\lecturerresponse{The principle that the laws of nature should have the same
form in every inertial frame was central to his reasoning. There was also
experimental evidence concerning light and the failure to detect motion
through an ether, although the lecturer is explicitly uncertain about which
results Einstein knew or how strongly they influenced him. The conceptual
pressure was that Maxwell's equations and the Lorentz force should obey the
relativity principle, which requires electric fields, magnetic fields, and
velocity components to transform together.}
\end{classroomqa}

Einstein's 1905 presentation did not use the compact four-vector notation
used here. It worked component by component, with the necessary factors of
\(\gamma\) and velocity written explicitly. Minkowski subsequently organized
the same structure geometrically as four-dimensional spacetime.

\begin{classroomqa}
\audiencequestion{Did Einstein resist Minkowski's spacetime formulation, and
how soon did Minkowski introduce it?}
\lecturerresponse{The recollection offered in class is that Einstein did not
initially welcome the four-dimensional geometric language but adopted its
value quickly. The timing is given tentatively as roughly two years after the
1905 paper, perhaps 1907; the speaker stresses that he is recalling rather
than establishing the history.}
\end{classroomqa}

\begin{classroomqa}
\audiencequestion{Was Minkowski a mathematician or a physicist, what else did
he contribute, and was he influenced by Lorentz or FitzGerald?}
\lecturerresponse{No confident answer is claimed. The distinction between
mathematician and physicist was not pursued, an audience correction interrupts
a broad remark about Minkowski's other work, and the possible lines of
influence are left expressly uncertain. The secure point for this discussion
is Minkowski's decisive organization of relativity in four-dimensional
geometric notation.}
\end{classroomqa}

\begin{classroomqa}
\audiencequestion{Were the transformations that mix \(\boldsymbol E\) and
\(\boldsymbol B\) already present in Einstein's 1905 work?}
\lecturerresponse{Yes, but written component by component rather than as one
field-tensor equation. The comparison made in class is with Maxwell's
component equations: the content can be present before the compact notation
that reveals its common structure. Minkowski's contribution was to make that
structure systematic in spacetime language.}
\end{classroomqa}

That returns us to the mathematical point with which we began. Components are
indispensable for calculation, but the law becomes intelligible when those
components are recognized as parts of one transforming object. The next task
is to use \(F_{\mu\nu}\) to write the charged particle's equation of motion in
a form that is the same in every inertial frame.
